This draft has several limitations.

First, the current strongest evidence is local and diagnostic.
The Qwen3-0.6B lane is modern enough to be useful, but it is not a full SOTA benchmark.
Larger models, more seeds, and standardized OpenUnlearning-style baselines would be needed for a broad claim.

Second, current stress pools are synthetic biography style.
They are controlled and reproducible, but future work should test more diverse domains and forgetting types.

Third, the HLC objective explored here may be too shallow.
K1/K2/K4 variants can improve short-horizon behavior, but prior locked results show that improvements do not always persist to $t=128$ or $t=512$.

Fourth, the current draft intentionally avoids external related-work citations.
The next writing pass should add exact citations for TOFU, NPO, machine unlearning, model editing, and durable/privacy-oriented unlearning benchmarks.

The next evidence gates are:

\begin{enumerate}[leftmargin=*]
  \item finish seeded Qwen3-0.6B long-512 HLC vs matched GA;
  \item add or confirm a third Qwen seed if compute allows;
  \item rerun final aggregate tables from canonical scripts;
  \item decide whether the results section is positive HLC evidence or a negative diagnostic study;
  \item add related-work citations and final figures.
\end{enumerate}
